\documentclass[Specialissue]{jsaiart}

\usepackage[ipaex]{pxchfon} % 日本語利用のため削除不可
\usepackage{profile-2e,graphicx} % 著者画像挿入のため削除不可

\jtitle{和文タイトル}
\jsubtitle{和文サブタイトル}

\author{%
 \name{著者1姓}{名}{Auther1 Roman Name}
 \affiliation{日本語所属名1}%
     {職種 in English e.g.) Engineer}%
     {user1@cyberagent.co.jp, http://cyberagent.co.jp/~user1/}
\and
 \name{著者2姓}{名}{Auther2 Roman Name}
 \affiliation{日本語所属名2}%
     {職種 in English e.g.) Researcher}%
     {user2@cyberagent.co.jp, http://cyberagent.co.jp/~user2/}
}

\begin{keyword}
キーワードとして2〜5個指定
\end{keyword}

\begin{summary}
200〜500 ワード
以内でsummaryを記す
\end{summary}

\begin{document}
\maketitle

\section{はじめに}
引用の例\cite{latex,texbook}
\section{従来の手法の問題点}
\section{今回の改良点}
\section{実験}
\subsection{実験手法}
\subsection{実験結果}
\subsubsection{実験1}
\subsubsection{実験2}
\section{まとめ}

\newpage
\begin{acknowledgment}
謝辞について
\end{acknowledgment}

\bibliography{btxsample}
\bibliographystyle{jsai}

\newpage
\appendix
\section{付録のタイトル1}
付録の本文1
\section{付録のタイトル2}
付録の本文2

% 著者の姓と名の間は半角スペースで区切る
% 略歴は200字以内
\begin{biography}
\profile{n}{著者1姓 名}{著者1の略歴}{portrait}
\end{biography}

\end{document}